%\documentclass[20pt,a1paper,portrait,margin=2mm, innermargin=18mm, blockverticalspace=15mm, colspace=15mm, subcolspace=8mm]{tikzposter} 

%\documentclass[20pt,a1paper,landscape,margin=2mm, innermargin=18mm, blockverticalspace=15mm, colspace=15mm, subcolspace=8mm]{tikzposter}

\documentclass[25pt,a0paper,portrait,margin=2mm, innermargin=30mm, blockverticalspace=20mm, colspace=20mm, subcolspace=10mm]{tikzposter}

%\documentclass[25pt,a0paper,landscape,margin=2mm, innermargin=28mm, blockverticalspace=20mm, colspace=20mm, subcolspace=10mm]{tikzposter}

\usepackage{aigpostersettings}

\usepackage{anyfontsize}
\usepackage{palatino}
\usepackage[utf8]{inputenc}
\usepackage[T1]{fontenc}
%\usepackage{cite}
\usepackage{booktabs}
\usepackage[table]{xcolor}
\usepackage{graphicx}
\usepackage{amssymb,amsthm,amsmath,amsfonts}

\usepackage[numbers]{natbib}
\renewcommand{\bibfont}{\footnotesize}
\bibliographystyle{abbrvnat}

\usepackage{serialisation}
%\tikzset{
%    serial node/.append={draw=black,line width=1pt,thick},
%}
\usepackage[macros,attackstyle=large,argumentstyle=colored,namestyle=bold]{argumentation}
\setargumentstyle{argument colored,font=\Large}
\setattackstyle{-{Stealth[scale=1.6]}}
\setargumentcolorscheme{aigblue}{aigyellow}

\theoremstyle{definition}
\newtheorem{definition}{Definition}
\newtheorem*{definition*}{Definition}

\usepackage{_macros}


\usetitlestyle{aigpostertitle1}
%%%%%%%%%%%%%%%%%%%%%%%%%%%%%%%%%%%%%%%%%%%%%%%

\title{\parbox{0.9\linewidth}{AgonProject}}
\institute{Artificial Intelligence Group, University of Hagen, Germany} 
%\author{Lars Bengel, Matthias Thimm}


%%%%%%%%%%%%%%%%%%%%%%%%%%%%%%%%%%%%%%%%%%%%%%%%%%%%%%%%%%%%%%

%Titlemaking

%Style-Inclusion from aigpostersettings.sty (check out title options for multiple logos)

\usecolorstyle{aigpostercolor}
\usebackgroundstyle{aigposterback}
\useblockstyle{aigposterblockB}

\begin{document}
	
\maketitle % See Section 4.1
	
%	\block{Abstract}{Blocktext} % See Section 4.2

\begin{columns}
\column{1.0}
\useblockstyle{aigposterblockY}
\block{Motivation}{
    \begin{minipage}{0.38\linewidth}
    \centering
     \begin{af}
            \argument{c}
            \argument[right=of a1]{o}
            \argument[right=of a2]{f}
            \argument[right=of a3]{s}
            \argument[right=of a4]{m}

            \dualattack{a1}{a2}
            \attack{a2}{a3}
            \attack{a3}{a4}
            \attack{a5}{a4}
        \end{af}

        \vspace*{2ex}

        \begin{minipage}{0.8\linewidth}
        \innerblock{}{
        \centering
        \textcolor{aigblue}{Goal:} Stepwise evaluation of the abstract dialectical 
        framework.
        }\end{minipage}
        \vspace{-0.3cm}
    \end{minipage}
    \hspace{5ex}
    \begin{minipage}{0.2\linewidth}
    \begin{description}
        \item[o:~] Valve is \textbf{o}pen.
        \item[c:~] Valve is \textbf{c}losed.
        \item[f:~] \textbf{F}low sensor activated.
        \item[s:~] System \textbf{s}hutdown.
        \item[m:~] System \textbf{m}aintenance.
    \end{description}
    \end{minipage}
    \hfill
    \begin{minipage}{0.34\linewidth}
    \begin{itemize}
        \item The main valve is either open or closed.
        \item If the valve is open (for too long), then the water flow sensor is activated.
        \item If the flow sensor is active or the system required maintenance, then the facility is shut down.
        \item The system requires no maintenance.
    \end{itemize}
    \end{minipage}
    \hspace{3ex}
}\useblockstyle{aigposterblockB}
\end{columns}


\begin{columns}
\column{0.5}

\block{Abstract Dialectical Frameworks}{
        \innerblock{}{
        \centering
        An \emph{abstract dialectical framework} (ADF) is a pair $\ADFcomplete$ where $\arguments$ is a set of arguments and $\conditions$ is a set of propositional formulae $\{\phi_a\}_{a\in\arguments}$ over $\arguments$, called \emph{acceptance conditions}. [Brewka and Woltran, 2010]
        }

        \vspace{1ex}
        Given an abstract dialectical framework \ADFcomplete.
        
        \begin{itemize}
            \item We consider three-valued interpretations: \modelcomplete
            \item We define the \emph{consensus} $\model_3 = \model_1 \sqcap \model_2$ as 
            \begin{align*}
                \model_3(a) = \begin{multicases}{3}
                    \tr & \text{if } \model_1(a) = \model_2(a) = \tr \\
                    \f & \text{if } \model_1(a) = \model_2(a) = \f\\
                    \un & \text{otherwise}         
                    \end{multicases}
            \end{align*}
            for all $a \in\arguments$.
    
            \item We define the set of \emph{completions} $[\model]_2$ as
            \begin{align*}
                [\model]_2 = \{\model' \mid \model \leq_i \model', (\model')^{-1}(\un) = \emptyset\}
            \end{align*}
    
            \item We define the \emph{characteristic operator} $\chop$ as
                \begin{align*}
                    \chop(\model)(a) = \bigsqcap\{\model'(\phi_a) \mid \model' \in [\model]_2\}
                \end{align*}
                for all $a\in\arguments$.
            \item An interpretation $\model: \arguments \mapsto \{\tr,\f,\un\}$ is called an \emph{admissible (\adm) model} of \ADF, iff $\model \leq_i \chop(\model)$.
        \end{itemize}
        \vspace*{-1.45cm}

        \node[](inf) at (38, 20.5){
        \begin{minipage}{0.5\linewidth}
        \begin{tikzpicture}
            \node[](u) at (0,0) {\un};
            \node[](t) at (-1.75,3.25) {\tr};
            \node[](f) at (1.75,3.25) {\f};
            \path(u) edge[->] node[left,yshift=-0.25cm]{$\leq_i$} (t);
            \path(u) edge[->] node[right,yshift=-0.2cm]{$\leq_i$} (f);
        \end{tikzpicture}
        \end{minipage}
        };
    
}

\column{0.5}

\block{Initial Models}{
    \begin{minipage}{.4\linewidth}
        % Innerblock Colors
	\colorlet{innerblocktitlebgcolor}{aigyellow}
	\colorlet{innerblocktitlefgcolor}{white}
	\colorlet{innerblockbodybgcolor}{aigyellow!30}
	\colorlet{innerblockbodyfgcolor}{black}
        \innerblock{Initial Model}{
        An interpretation $\model$ is called an \emph{initial model} (\elc) of $\ADF$, iff $\model$ is admissible with $\model \neq \model_{\un}$ and there is no admissible model $\model'\neq \model_{\un}$ with $\model' <_i \model$.
        }
    \end{minipage}%
    \hfill
    \begin{minipage}{0.57\linewidth}
        \vspace*{0.4cm}
        An initial model $\model$ is
	\begin{itemize}
        \item \emph{unattacked} iff for all $a\in\arguments$ if $\model(a)\neq\un$, then $\phi_a \in \textsc{Taut}_\arguments\cup\textsc{Unsat}_\arguments$,
		\item \emph{unchallenged} iff \model is not unattacked and there is no initial model $\model'$ that is conflicting with $\model$,
		\item \emph{challenged} iff there is some model initial model $\model'$ that is conflicting with $\model$.
	\end{itemize}
        
    \end{minipage}
}

\useblockstyle{aigposterblockY}
\block{Example: Initial Models}{
    \begin{center}
        \begin{af}
        \argument{c}
        \argument[right=of a1]{o}
        \argument[right=of a2]{f}
        \argument[right=of a3]{s}
        \argument[right=of a4]{m}

        \dualattack{a1}{a2}
        \attack{a2}{a3}
        \attack{a3}{a4}
        \attack{a5}{a4}
        \end{af}
    \end{center}

    \centering
    \vspace{1cm}

    \highlight{Initial Models}\\[0.7ex]    
    \begin{tabular}{clclc}
        unattacked && unchallenged && challenged \\
        $\{m \mapsto \f\}$ &\phantom{xx}& -- &\phantom{xx}& $\{c \mapsto\tr, o\mapsto\f\}$, $\{c\mapsto\f,o\mapsto\tr\}$
    \end{tabular}
    \vspace{-0.175cm}
}\useblockstyle{aigposterblockB}

\end{columns}
\begin{columns}

\column{0.4}
%\useblockstyle{}
\block{\model-Reduct}{
    \vspace{0.3cm}
    Let \ADFcomplete be an ADF and $\modelcomplete$ is a three-valued interpretation.
    Then we define the $\model$\emph{-reduct} of $\ADF$ as the ADF $\ADF^\model = (\arguments', \{\phi'_a\}_{a\in\arguments'})$ where
    \begin{align*}
        \arguments' = &\quad \arguments \setminus \{a\in\arguments\mid \model(a)\neq\un\},\\
        \conditions' = &\quad \{\phi'_a\}_{a\in\arguments'},
    \end{align*}
    with $\mathsf{x}\in\{\tr, \f\}$ and $\phi'_a = \phi_a^{[b/\mathsf{x}~:~\model(b)=\mathsf{x}]}$.
    \vspace{-0.1cm}
}

\column{0.6}

\block{Serialisation Sequences}{
    \vspace{0.5cm}
    \begin{minipage}{.4\linewidth}
        \colorlet{innerblocktitlebgcolor}{aigyellow}
	\colorlet{innerblocktitlefgcolor}{white}
	\colorlet{innerblockbodybgcolor}{aigyellow!30}
	\colorlet{innerblockbodyfgcolor}{black}
        \innerblock{Serialisation Sequence}{
            A \emph{serialisation sequence} for \ADFcomplete is a sequence \msequencecomplete with $\model_1 \in \elc(\ADF)$ and for each $2 \leq i \leq n$ we have that $\model_i \in \elc(\ADF^{\model_1 \sqcup\dots\sqcup\model_{i-1}})$.
        }
    \end{minipage}
    \hfill
    \begin{minipage}{0.56\linewidth}
        \begin{itemize}
            \item Provides a decomposition of admissible models into a series of initial models.
            \item Each sequence represents one way of reasoning to arrive at some conclusion.
            \item Directly encodes procedural information in the semantic representation
        \end{itemize}
    \end{minipage}
    \vspace{-0.2cm}
}
\end{columns}

\begin{columns}
    \column{1}
    \useblockstyle{aigposterblockY}
    \block{Example: Serialisation Sequences}{
    \begin{minipage}{0.30\linewidth}
        \centering
        \begin{af}
            \argument{c}
            \argument[right=of a1]{o}
            \argument[right=of a2]{f}
            \argument[right=of a3]{s}
            \argument[right=of a4]{m}

            \annotation[yshift=-0.8cm]{a1}{$\phi_c = \neg o$}
            \annotation[yshift=-0.8cm]{a2}{$\phi_o = \neg c$}
            \annotation[yshift=-0.8cm]{a3}{$\phi_f = o$}
            %\annotation<3,4>[yshift=-0.8cm]{a3}{$\phi_f = \bot$}
            \annotation[yshift=-0.8cm]{a4}{$\phi_s = f \lor m$}
            %\annotation<5,6>[yshift=-0.8cm]{a4}{$\phi_s = \bot \lor m$}
            %\annotation<7,8>[yshift=-0.8cm]{a4}{$\phi_s = \bot$}
            \annotation[yshift=-0.8cm]{a5}{$\phi_m = \bot$}
    
            \dualattack{a1}{a2}
            \attack{a2}{a3}
            \attack{a3}{a4}
            \attack{a5}{a4}
        \end{af}

        \vspace{0.5cm}

        Initial Models:
        
        $\mathhighlight{\{c \mapsto \tr, o \mapsto \f\}}, \{c \mapsto \f, o \mapsto \tr\}, \{m \mapsto \f\}$
        
            
        
    \end{minipage}
    \hfill
    \begin{minipage}{0.25\linewidth}
        \centering
        \begin{af}
            %\argument{c}
            %\argument[right=of a1]{o}
            \argument(a3){f}
            \argument[right=of a3](a4){s}
            \argument[right=of a4](a5){m}

            %\argument(x){x} at (20,0)
            \afname{$\implies$} at (-5,0)
            \afname{$\ADF^\model$} at (-4.9,1)

            \annotation[yshift=-0.8cm]{a3}{$\phi_f = \bot$}
            \annotation[yshift=-0.8cm]{a4}{$\phi_s = f \lor m$}
            %\annotation<5,6>[yshift=-0.8cm]{a4}{$\phi_s = \bot \lor m$}
            %\annotation<7,8>[yshift=-0.8cm]{a4}{$\phi_s = \bot$}
            \annotation[yshift=-0.8cm]{a5}{$\phi_m = \bot$}
    
            \attack{a3}{a4}
            \attack{a5}{a4}
        \end{af}

        \vspace{0.5cm}

        Initial Models:
        
        $\mathhighlight{\{f \mapsto \f\}}, \{m \mapsto \f\}$
    \end{minipage}
    \hfill
    \begin{minipage}{0.16\linewidth}
        \centering
        \begin{af}
            %\argument{c}
            %\argument[right=of a1]{o}
            %\argument[right=of a2]{f}
            \argument(a4){s}
            \argument[right=of a4](a5){m}

            \afname{$\implies$} at (-4,0)
            \afname{$\ADF^\model$} at (-3.9,1)

            \annotation[yshift=-0.8cm]{a4}{$\phi_s = \bot \lor m$}
            %\annotation<7,8>[yshift=-0.8cm]{a4}{$\phi_s = \bot$}
            \annotation[yshift=-0.8cm]{a5}{$\phi_m = \bot$}
    
            \attack{a5}{a4}
        \end{af}

        \vspace{0.5cm}

        Initial Models:
        
        $\mathhighlight{\{m \mapsto \f\}}$
        
    \end{minipage}
    \hfill
    \begin{minipage}{0.10\linewidth}
        \centering
        \begin{af}
            \argument(a4){s}

            \afname{$\implies$} at (-5,0)
            \afname{$\ADF^\model$} at (-4.9,1)

            \annotation[yshift=-0.8cm]{a4}{$\phi_s = \bot$}
        \end{af}

        \vspace{0.5cm}

        Initial Models:
        
        $\mathhighlight{\{s \mapsto \f\}}$
    \end{minipage}

    \centering
    \vspace{0.9cm}
    
    Serialisation sequence: $\mseq = (\{c \mapsto \tr, o \mapsto \f\}, \{f \mapsto \f\}, \{m \mapsto \f\}, \{s \mapsto \f\})$, ~~corresponding to the admissible model $\{c \mapsto \tr, o \mapsto \f, f \mapsto \f, s \mapsto \f, m \mapsto \f\}$

    \vspace{-0.3cm}
    
    }\useblockstyle{aigposterblockB}
\end{columns}

\begin{columns}
    \column{0.6}
    \block{Characterising Admissibility-based Semantics}{
    We can characterise many semantics in terms of serialisation sequences as follows.\\[0.4ex]

        \hspace{1.8cm}
        \begin{minipage}{0.76\linewidth}
        \msequencecomplete is a serialisation sequence of \ADF and \dots
        \begin{description}
            \item[preferred:] if there exists no further initial model in the final reduct.
            \item[complete:] if there exists no unattacked initial model in the final reduct.
            \item[grounded:] if every $\model_i$ is an unattacked initial model and there exists no further unattacked initial model.
            \item[two-valued:] if the final reduct is the empty ADF.
        \end{description}
        \end{minipage}

        \centering

        \vspace{1cm}

        \yellowhighlight{
        $\implies$ Serialisation Sequences provide a more expressive semantic representation for ADF.
        }
    }
    \column{0.4}
    \block{Complexity of Initial Models}{
        \renewcommand{\arraystretch}{1.25}
        \centering
        %\begin{minipage}{\textwidth}
        %\centering
        \begin{tabular}{l|cc|cc}
        \toprule
        & \multicolumn{2}{c|}{AFs} & \multicolumn{2}{c}{ADFs}\\
         & \adm [1] & \elc [2] & \adm [3] &\cellcolor{aigyellow!10} \elc \\
        \midrule
        $\textsc{Ver}_{\sigma}$ & in \textsf{L} & in \p & \coNP-c &\cellcolor{aigyellow!10} in $\p^{\NP}_\parallel$, \coNP-h \\
        $\textsc{Exists}_{\sigma}$ & trivial & \NP-c & trivial &\cellcolor{aigyellow!10} $\Sigma^P_2$-c \\
        $\textsc{Exists}^{\neg\emptyset}_{\sigma}$ & \NP-c & \NP-c & $\Sigma^P_2$-c &\cellcolor{aigyellow!10} $\Sigma^P_2$-c \\
        $\textsc{Cred}_{\sigma}$ & \NP-c & \NP-c & $\Sigma^P_2$-c &\cellcolor{aigyellow!10} $\Sigma^P_2$-c \\
        $\textsc{Skept}_{\sigma}$ & trivial & \coNP-c & trivial &\cellcolor{aigyellow!10} $\Pi^P_2$-c \\
        \bottomrule
        \end{tabular}
        %\end{minipage}

        \vspace{0.4cm}
        {\small
        [1] Dvorak and Dunne, 2017; [2] Thimm, 2022; [3] Strass and Wallner, 2015
        }\vspace{-0.2cm}
    }
\end{columns}

\node[] at (38.3,43.9){\bf\ttfamily\LARGE \textcolor{white}{1}};
\node[] at (-1.5,29.2){\bf\ttfamily\LARGE \textcolor{white}{2}};
\node[] at (38.3,29.2){\bf\ttfamily\LARGE \textcolor{white}{3}};
\node[] at (38.3,12){\bf\ttfamily\LARGE \textcolor{white}{4}};
\node[] at (-9.5,-5){\bf\ttfamily\LARGE \textcolor{white}{5}};
\node[] at (38.3,-5){\bf\ttfamily\LARGE \textcolor{white}{6}};
\node[] at (38.3,-19.5){\bf\ttfamily\LARGE \textcolor{white}{7}};
\node[] at (6.4,-36.6){\bf\ttfamily\LARGE \textcolor{white}{8}};
\node[] at (38.3,-36.6){\bf\ttfamily\LARGE \textcolor{white}{9}};

%\hfill
\end{document}